\chapter{Public API and Design Contract}

\section{Opaque Heap Handle}

Client code interacts with heapx through an opaque \texttt{struct heapx\_heap}.
The concrete backend layout is private.
This prevents callers from depending on implementation details and allows the same public functions to dispatch to different heap implementations.

\begin{lstlisting}[style=cstyle, caption={Creating a heap through the public API.}]
struct heapx_heap *heapx_create(
    enum heapx_implementation implementation,
    heapx_cmp_fn cmp
);
\end{lstlisting}

The backend is selected at construction time through \texttt{enum heapx\_implementation}.
After creation, client code uses the same operations regardless of whether the heap is backed by a binary heap, a Fibonacci heap, or the Kaplan heap.

\section{Comparator Contract}

heapx stores generic \texttt{void *} items.
Ordering is defined by a caller-provided comparator:

\begin{lstlisting}[style=cstyle, caption={Comparator type.}]
typedef int (*heapx_cmp_fn)(const void *lhs, const void *rhs);
\end{lstlisting}

A negative result means that \texttt{lhs} has higher priority and should appear closer to the minimum.
The library therefore implements min-priority behavior relative to the comparator, not relative to a fixed numeric key.

The comparator must remain valid for the lifetime of the heap.
It should define a consistent ordering for all objects stored in the same heap.
If a stored object is mutated so that its priority improves, the caller must use \texttt{heapx\_decrease\_key()} with the item's handle.
If the object is mutated in the opposite direction, heapx does not provide a repair operation for that case.

\section{Ownership}

The library never owns stored items.
It stores item pointers exactly as provided by the caller:

\begin{itemize}
    \item heapx does not copy item objects;
    \item heapx does not free item objects;
    \item destroying a heap releases only heap-owned metadata;
    \item callers must keep stored objects valid until they are removed or the heap is destroyed.
\end{itemize}

This design is idiomatic for a small C library and keeps heapx independent from item lifetime policies.
It also makes the allocation behavior of heap metadata easier to reason about during benchmarks.

\section{Core Operations}

The public operations are intentionally narrow:

\begin{lstlisting}[style=cstyle, caption={Core heap operations.}]
int heapx_insert(struct heapx_heap *heap, void *item);
int heapx_insert_handle(
    struct heapx_heap *heap,
    void *item,
    struct heapx_handle *out
);
int heapx_decrease_key(
    struct heapx_heap *heap,
    struct heapx_handle handle
);
void *heapx_remove(
    struct heapx_heap *heap,
    struct heapx_handle handle
);
void *heapx_peek_min(const struct heapx_heap *heap);
void *heapx_extract_min(struct heapx_heap *heap);
size_t heapx_size(const struct heapx_heap *heap);
int heapx_empty(const struct heapx_heap *heap);
\end{lstlisting}

\texttt{heapx\_insert()} is suitable when the caller only needs normal priority queue behavior.
\texttt{heapx\_insert\_handle()} additionally returns a handle that can later be used for targeted operations.

\section{Targeted Operations}

\texttt{heapx\_decrease\_key()} and \texttt{heapx\_remove()} operate on handles, not by searching for an item pointer.
This is important for preserving the intended complexity of heap operations.
Searching for an arbitrary pointer would be linear and would obscure the algorithmic behavior of the backend.

The handle mechanism is discussed in detail in Chapter~\ref{chap:handles}.
At the API level, the important contract is simple:

\begin{itemize}
    \item a handle is live only while the associated item remains stored in the heap that created it;
    \item after extraction or removal, the handle becomes stale;
    \item stale handles fail cleanly;
    \item handles from another heap fail cleanly.
\end{itemize}

\section{Defined Edge Cases}

The public dispatch layer centralizes common edge cases.
\texttt{heapx\_destroy(NULL)} does nothing.
\texttt{heapx\_size(NULL)} returns zero.
\texttt{heapx\_empty(NULL)} treats a null heap as empty.
\texttt{peek\_min} and \texttt{extract\_min} return \texttt{NULL} for a null or empty heap.
Insert and decrease-key operations return \texttt{-1} when the heap pointer or handle is invalid.

These rules make backend implementations simpler because public null-handling is not duplicated in every backend.

\section{The Role of contains}

\texttt{heapx\_contains()} checks pointer identity.
It is useful as a convenience query and as a diagnostic operation in tests, but it is not a heap-native performance primitive.
The current backends implement it as a linear search.

For this reason, complexity tables and benchmark interpretation should focus on insert, insert-handle, decrease-key, remove, peek-min, extract-min, size, and empty.
